% GENERATED by build_chapters.py; edit docs/report/chapters, not this file

\chapter{Security, Testing and Deployment}

\section{Security}

A pipeline that reads logs written by anyone who can push code, sends them to a model, and then writes to a repository has three things to protect: its credentials, the repository, and its own decisions. This chapter lists what is protected and how, and ends with the weaknesses that are known and open.

\subsection{Credentials}

\begingroup\small
\begin{longtable}{>{\raggedright\arraybackslash}p{0.165\dimexpr\linewidth-6\tabcolsep\relax}>{\raggedright\arraybackslash}p{0.457\dimexpr\linewidth-6\tabcolsep\relax}>{\raggedright\arraybackslash}p{0.378\dimexpr\linewidth-6\tabcolsep\relax}}
\toprule
\textbf{secret} & \textbf{where it lives} & \textbf{who can read it} \\
\midrule
\endhead
\bottomrule
\endfoot
\endlastfoot
Webhook shared secret & the n8n credential "N8N Webhook Secret" (header auth, header \texttt{X-\allowbreak{}Webhook-\allowbreak{}Secret}), the GitHub repository secret \texttt{N8N\_\allowbreak{}WEBHOOK\_\allowbreak{}SECRET}, and \texttt{secrets.\allowbreak{}json} on the operator's PC & the workflow, the notify job and the launcher scripts \\
GitHub token & the n8n credential "GitHub PAT" (bearer); locally taken from the GitHub CLI by \texttt{tools/\allowbreak{}golive/\allowbreak{}n8n\_\allowbreak{}import.\allowbreak{}py} & the n8n HTTP nodes that call the GitHub API \\
Model endpoint key & the n8n credential "Qwen3.8 Endpoint Key" and the constant \texttt{API\_\allowbreak{}KEY} in \texttt{kaggle/\allowbreak{}endpoint/\allowbreak{}deploy\_\allowbreak{}src.\allowbreak{}py} & anyone who can read the repository \\
Webhook address & the repository secret \texttt{N8N\_\allowbreak{}WEBHOOK\_\allowbreak{}URL} & the notify job \\
\bottomrule
\caption{Credentials}\label{tab:10securi-1}\\
\end{longtable}
\endgroup

Runtime secrets (\texttt{secrets.\allowbreak{}json}, the saved model address) live under \texttt{\%LOCALAPPDATA\%\textbackslash{}selfheal-\allowbreak{}n8n\textbackslash{}}, outside the repository. The backend's \texttt{.\allowbreak{}env} is ignored by git. The repository contains no token strings; this is checked by \texttt{tools/\allowbreak{}tests/\allowbreak{}test\_\allowbreak{}workflow\_\allowbreak{}structure.\allowbreak{}py} for the workflow JSON, and by a search of the whole tree before each push (see below). Credentials are imported into a local n8n as a plaintext file that the import script deletes immediately afterwards (tools/\allowbreak{}golive/\allowbreak{}n8n\_\allowbreak{}import.py).

\subsection{The webhooks}

Both n8n webhooks require the shared secret as a header and n8n rejects a request without it before the workflow starts: a wrong secret gave HTTP 403 in a test against the cloud instance, and the right one started the workflow. The second webhook, which stores the model address, accepts only an address that matches \texttt{https:\allowbreak{}/\allowbreak{}/\allowbreak{}\textless{}letters, digits, hyphens\textgreater{}.\allowbreak{}trycloudflare.\allowbreak{}com} and keeps the old address otherwise (tools/\allowbreak{}nodes/\allowbreak{}store\_\allowbreak{}model\_\allowbreak{}url.js, tested in tools/\allowbreak{}tests/\allowbreak{}test\_\allowbreak{}model\_\allowbreak{}url.py with look-alike hosts such as \texttt{https:\allowbreak{}/\allowbreak{}/\allowbreak{}abc.\allowbreak{}trycloudflare.\allowbreak{}com.\allowbreak{}evil.\allowbreak{}com}). A stolen secret could therefore still point the workflow at an attacker's tunnel, which is one reason to keep the secret long and to rotate it.

\texttt{ci-\allowbreak{}heal-\allowbreak{}backend} verifies GitHub's HMAC-SHA256 signature on the raw body with a constant-time comparison, applies a repository allowlist that denies by default, protects its operator API with a token outside development, and now ignores runs from fork pull requests and refuses to start outside \texttt{dev} with an empty or placeholder secret (docs/\allowbreak{}results/\allowbreak{}BACKEND\_\allowbreak{}AUDIT.md).

\subsection{The log redirect}

The job-log endpoint of the GitHub API answers with a redirect to a signed address on another host. n8n re-sends the credential's authorization header when it follows a redirect, which would hand the GitHub token to the storage host. The workflow therefore does not follow the redirect: one node reads the address from the 302 answer, and a second node without any credential downloads it. The real-engine rig checks this: the token reaches the API host and does not reach the blob host (tools/\allowbreak{}n8n\_\allowbreak{}rig/\allowbreak{}drive\_\allowbreak{}full.py, eval/\allowbreak{}runs/\allowbreak{}testall/\allowbreak{}n8n\_\allowbreak{}rig.txt).

\subsection{Treating the model as untrusted}

The prompt contains text written by whoever caused the failure: log lines, file contents. The design assumes the model may be steered by that text or may simply be wrong, so every step after the reply is a gate.

\begin{itemize}[leftmargin=*, itemsep=2pt]
\item Only failures marked \texttt{failure\_\allowbreak{}kind=code} reach the model, and a payload with any other kind is refused inside the workflow as well (tools/\allowbreak{}nodes/\allowbreak{}build\_\allowbreak{}ai\_\allowbreak{}prompt.js).
\item The reply must be a JSON object with SEARCH/\allowbreak{}REPLACE edits. Each SEARCH text must occur in the file exactly once, and the result must parse as JavaScript (tools/\allowbreak{}nodes/\allowbreak{}extract\_\allowbreak{}ai\_\allowbreak{}response.js).
\item The file the edit names must be one of the files the model was shown.
\item A branch named \texttt{main} or \texttt{master} is replaced by an \texttt{ai-\allowbreak{}fix-\allowbreak{}} name, and the fix is a pull request, never a push to the main branch.
\item CI runs on the pull request, and a person merges it.
\item The backend adds secret redaction before any text reaches a model, caps on edit size, a check against weakened tests, and an \texttt{ai-\allowbreak{}fix-\allowbreak{}} prefix rule enforced in the GitHub client (ci-heal-backend/\allowbreak{}app/\allowbreak{}services/\allowbreak{}github\_\allowbreak{}client.py, \texttt{\_\allowbreak{}refuse}).
\end{itemize}

\subsection{Known weaknesses}

\begin{itemize}[leftmargin=*, itemsep=2pt]
\item \textbf{The endpoint key is public.} It is a constant in a file in the repository. Anyone who reads the repository and learns the tunnel address can use the GPU session. The address changes on every boot and is not published, which limits the exposure but is not authentication.
\item \textbf{The shared secret in use is short.} It was chosen for convenience and should be replaced by a long random value before the repository or its secrets are shared.
\item \textbf{A misled fix is still a candidate.} A log line that argues for a particular edit can influence the model. The gates above stop a malformed or unrelated edit, not a plausible wrong one that happens to pass the tests. The review of the pull request is the control.
\item \textbf{Pull-request CI runs the proposed code.} The fix pull request is built by the repository's own workflow, with whatever that workflow can access. The audit of the backend noted that only workflow files, \texttt{.\allowbreak{}git} and \texttt{.\allowbreak{}env} files are blocked from editing, so a build entry point such as \texttt{package.\allowbreak{}json} can be changed by a fix.
\item \textbf{Seven workflow nodes have no error setting}, so a failure there ends the run silently (Chapter 11). Seventeen backend findings are open (docs/\allowbreak{}results/\allowbreak{}BACKEND\_\allowbreak{}AUDIT.md).
\item \textbf{The backend reads a webhook body in full} before it checks the size limit and the signature (finding 11 in docs/\allowbreak{}results/\allowbreak{}BACKEND\_\allowbreak{}AUDIT.md).
\end{itemize}

\subsection{Checks before publishing}

Before each push the tree is searched for the token prefixes \texttt{ghp\_\allowbreak{}} and \texttt{github\_\allowbreak{}pat\_\allowbreak{}}, for \texttt{hf\_\allowbreak{}}, for the webhook secret value, and for \texttt{.\allowbreak{}env} and \texttt{secrets.\allowbreak{}json} files. A match in a test that deliberately contains a fake token is accepted after reading it.

\section{Testing and Verification}

The project has four layers of checks, from fast and local to slow and live. The first three need no GPU and write nothing to GitHub. A fuller guide with commands is in docs/\allowbreak{}guides/\allowbreak{}testing.md.

\subsection{Layer 1: backend unit tests}

\texttt{ci-\allowbreak{}heal-\allowbreak{}backend} has a pytest suite that uses a mock model and a mocked GitHub client, so it needs no key and no network. It collected 157 tests before the audit fixes described in Chapter 8 and 205 after, all passing, with \texttt{ruff check .\allowbreak{}} and \texttt{ruff format -\allowbreak{}-\allowbreak{}check .\allowbreak{}} clean (run on 2026-11-15, output recorded in docs/\allowbreak{}results/\allowbreak{}BACKEND\_\allowbreak{}AUDIT.md). The suite covers the classifier, the decision rules, patch application, the HTTP model client, the GitHub client, both webhooks, secret redaction, log trimming, the database layer and the settings.

\subsection{Layer 2: node suites and eval scorers}

\texttt{python run\_\allowbreak{}tests.\allowbreak{}py} runs every suite in \texttt{tools/\allowbreak{}tests/\allowbreak{}}, the scorer self-tests in \texttt{eval/\allowbreak{}} and a generator check. Each Code node of the n8n workflow is a plain JavaScript file, so \texttt{tools/\allowbreak{}heal\_\allowbreak{}harness.\allowbreak{}py} can run it under Node with a fake n8n context. The table gives the number of passing assertions each suite printed in its last run (eval/\allowbreak{}runs/\allowbreak{}testall/\allowbreak{}).

\begingroup\small
\begin{longtable}{>{\raggedright\arraybackslash}p{0.397\dimexpr\linewidth-6\tabcolsep\relax}>{\raggedright\arraybackslash}p{0.190\dimexpr\linewidth-6\tabcolsep\relax}>{\raggedright\arraybackslash}p{0.414\dimexpr\linewidth-6\tabcolsep\relax}}
\toprule
\textbf{suite} & \textbf{passing assertions} & \textbf{what it checks} \\
\midrule
\endhead
\bottomrule
\endfoot
\endlastfoot
test\_\allowbreak{}build\_\allowbreak{}prompt & 29 & the prompt carries the trimmed CI error and the shipped system prompt \\
test\_\allowbreak{}extract\_\allowbreak{}node & 15 & SEARCH/\allowbreak{}REPLACE application and every rejection reason \\
test\_\allowbreak{}assemble\_\allowbreak{}reply & 15 & rebuilding a streamed reply, and the failure explanation \\
test\_\allowbreak{}heal\_\allowbreak{}outcome & 6 & the retry outcome against GitHub error bodies \\
test\_\allowbreak{}model\_\allowbreak{}url & 30 & storing and reading the model address, and the fallback reply \\
test\_\allowbreak{}multifile & 19 & candidate file choice and prompt for multi-file bugs \\
test\_\allowbreak{}workflow\_\allowbreak{}structure & 22 & wiring, webhooks and credentials of the shipped workflow JSON \\
test\_\allowbreak{}shipped\_\allowbreak{}eval & 49 & the scorer of the quality study, including 15 of 15 real captured replies \\
check\_\allowbreak{}multifile\_\allowbreak{}cases & 7 & every multi-file case is a real test: golden passes, broken fails \\
test\_\allowbreak{}multifile\_\allowbreak{}eval & 13 & scoring of the multi-file study \\
test\_\allowbreak{}bakeoff\_\allowbreak{}eval & 52 & the bake-off kernel's functions and decision rule \\
\bottomrule
\caption{Layer 2: node suites and eval scorers}\label{tab:10testin-1}\\
\end{longtable}
\endgroup

The generator check regenerates \texttt{n8n/\allowbreak{}self\_\allowbreak{}healing\_\allowbreak{}pipeline.\allowbreak{}json} from \texttt{tools/\allowbreak{}gen\_\allowbreak{}n8n.\allowbreak{}py} and the node files and requires a byte-identical result, so the file in the repository cannot drift from its source. In the last full run all 13 checks passed (12 without the rig, 13 with it).

The harness has a known blind spot, stated in its own docstring: it resolves \texttt{\$('Node')} lookups that n8n's sandbox does not resolve for Code nodes. A passing harness run therefore does not prove that an accessor works in n8n. This is the reason for layer 3.

\subsection{Layer 3: a real n8n engine}

\texttt{python run\_\allowbreak{}tests.\allowbreak{}py -\allowbreak{}-\allowbreak{}rig} imports the workflow nodes into a real n8n 1.x and runs them against a local mock of GitHub and the model (tools/\allowbreak{}n8n\_\allowbreak{}rig/\allowbreak{}). The nodes under test are the shipped ones, copied from the workflow JSON; only host names are changed and credentials removed. The last run printed 18 passing scenarios and no failures (eval/\allowbreak{}runs/\allowbreak{}testall/\allowbreak{}n8n\_\allowbreak{}rig.txt). They include:

\begin{itemize}[leftmargin=*, itemsep=2pt]
\item the CI log reaches the prompt and the PAT is sent to the GitHub API host but not to the blob host the log redirect points to;
\item the same fix is produced on a first heal, on a retry (thinking on, same branch, trailer \texttt{Heal-\allowbreak{}Attempt:\allowbreak{} 2}) and on a trace-only multi-file bug;
\item each rejection path of the extract node posts a comment on the commit that quotes its reason;
\item a model call that times out with HTTP 524 and a fallback that also fails ends in a pull-request comment that names both causes;
\item a model call that times out and a fallback that answers still produces the correct fix and no comment;
\item a job log that returns 404 still heals with the text \texttt{(log unavailable)};
\item a payload with hostile characters still builds a valid request;
\item a payload with \texttt{failure\_\allowbreak{}kind=infra} is refused before any model call.
\end{itemize}

\subsection{Layer 4: live runs}

Live runs push a deliberately broken commit to a GitHub repository and watch the whole chain, with the real model on Kaggle. They are run by hand because they write to GitHub and use GPU quota. Chapter 9 and Chapter 10 report them.

\subsection{Audit}

A read-only security and reliability review of \texttt{ci-\allowbreak{}heal-\allowbreak{}backend} listed 21 findings. Four were fixed with tests that failed first, and 17 are open (docs/\allowbreak{}results/\allowbreak{}BACKEND\_\allowbreak{}AUDIT.md). The n8n workflow was reviewed by reading its nodes against the tests: seven nodes (Fetch Job Logs, Fetch Changed Files, Build AI Prompt, Get Model URL, Get File SHA, Push Fix File and Post Failure Comment) have no error setting, so a failure there ends the execution without a comment on GitHub (docs/\allowbreak{}architecture/\allowbreak{}n8n-workflow.md).

\subsection{What is not tested}

\begin{itemize}[leftmargin=*, itemsep=2pt]
\item \texttt{kaggle/\allowbreak{}endpoint/\allowbreak{}deploy\_\allowbreak{}src.\allowbreak{}py} runs only on a Kaggle GPU; its test is a real boot.
\item The GitHub write steps of the workflow, from branch creation to the pull request, run only in live runs.
\item The backend with \texttt{DRY\_\allowbreak{}RUN=false} against real GitHub has no automated test.
\item The OpenAI fallback node is tested for wiring, for shipping without a credential and for its error handling, but never against a real OpenAI account. On the n8n Cloud trial the included free credits answered \texttt{Forbidden} to a webhook-triggered run in our test (docs/\allowbreak{}guides/\allowbreak{}cloud-n8n.md).
\item \texttt{demo-\allowbreak{}ci.\allowbreak{}yml} is exercised only by pushing to GitHub.
\item The launcher scripts have no tests.
\end{itemize}

\section{Deployment and Operation}

The launcher folder holds double-click scripts for every step, so nobody has to remember the command lines (launcher/\allowbreak{}README.md). This chapter lists what runs where, the order to start things in, and what each step costs.

\subsection{What runs where}

\begingroup\small
\begin{longtable}{>{\raggedright\arraybackslash}p{0.163\dimexpr\linewidth-8\tabcolsep\relax}>{\raggedright\arraybackslash}p{0.153\dimexpr\linewidth-8\tabcolsep\relax}>{\raggedright\arraybackslash}p{0.432\dimexpr\linewidth-8\tabcolsep\relax}>{\raggedright\arraybackslash}p{0.253\dimexpr\linewidth-8\tabcolsep\relax}}
\toprule
\textbf{part} & \textbf{runs on} & \textbf{started by} & \textbf{stopped by} \\
\midrule
\endhead
\bottomrule
\endfoot
\endlastfoot
Qwen3.8-27B endpoint & Kaggle, 2 x T4 & \texttt{launcher/\allowbreak{}start/\allowbreak{}1\_\allowbreak{}start\_\allowbreak{}model\_\allowbreak{}on\_\allowbreak{}kaggle.\allowbreak{}bat} & Stop Session in the Kaggle page (the Kaggle command line cannot stop a session) \\
n8n workflow, cloud copy & n8n Cloud trial, always on & already published & unpublish in the editor \\
n8n workflow, local copy & this PC, \texttt{localhost:\allowbreak{}5678} & \texttt{start/\allowbreak{}start\_\allowbreak{}n8n\_\allowbreak{}local\_\allowbreak{}only.\allowbreak{}bat}, or \texttt{2\_\allowbreak{}} and \texttt{3\_\allowbreak{}} for the live setup & \texttt{end/\allowbreak{}stop\_\allowbreak{}n8n.\allowbreak{}bat} \\
ci-heal-backend & this PC, port 8000 & \texttt{start/\allowbreak{}4\_\allowbreak{}start\_\allowbreak{}backend.\allowbreak{}bat} & \texttt{end/\allowbreak{}stop\_\allowbreak{}backend.\allowbreak{}bat} \\
demo repository workflow & GitHub Actions & a push to \texttt{main} that changes \texttt{demo-\allowbreak{}app/\allowbreak{}} & not applicable \\
\bottomrule
\caption{What runs where}\label{tab:10deploy-1}\\
\end{longtable}
\endgroup

\subsection{The live demo in order}

\begin{enumerate}[leftmargin=*, itemsep=2pt]
\item \texttt{1\_\allowbreak{}start\_\allowbreak{}model\_\allowbreak{}on\_\allowbreak{}kaggle.\allowbreak{}bat}. It prints the status of the project's kernels, builds the endpoint notebook from \texttt{kaggle/\allowbreak{}endpoint/\allowbreak{}deploy\_\allowbreak{}src.\allowbreak{}py}, pushes it, follows the boot log for about four minutes, saves the tunnel address to \texttt{\%LOCALAPPDATA\%\textbackslash{}selfheal-\allowbreak{}n8n\textbackslash{}qwen\_\allowbreak{}url.\allowbreak{}txt} and posts it to the cloud workflow's Set Model URL webhook.
\item Push a commit that breaks a test in \texttt{demo-\allowbreak{}app/\allowbreak{}}. The Actions run fails, the notify job posts to the n8n webhook, and a pull request on an \texttt{ai-\allowbreak{}fix-\allowbreak{}.\allowbreak{}.\allowbreak{}.\allowbreak{}} branch appears about a minute later.
\item Stop the Kaggle session when finished, either with \texttt{end/\allowbreak{}stop\_\allowbreak{}model\_\allowbreak{}on\_\allowbreak{}kaggle.\allowbreak{}bat} (it opens the notebook page and then checks that the session stopped) or directly in the page: View Active Events, the three-dot menu on the session, Stop Session.
\end{enumerate}

For the local copy of n8n instead of the cloud one, run \texttt{2\_\allowbreak{}load\_\allowbreak{}workflow\_\allowbreak{}into\_\allowbreak{}n8n.\allowbreak{}bat} and \texttt{3\_\allowbreak{}start\_\allowbreak{}n8n\_\allowbreak{}live.\allowbreak{}bat} after step 1. The second imports the workflow and credentials with the current model address, and the third starts n8n with a public Cloudflare tunnel and points the repository's \texttt{N8N\_\allowbreak{}WEBHOOK\_\allowbreak{}URL} secret at it.

\subsection{Cost and limits}

The Kaggle session uses GPU quota for as long as it runs. The kernel stops itself after 1.5 hours (kaggle/\allowbreak{}endpoint/\allowbreak{}deploy\_\allowbreak{}src.py, \texttt{HOURS}), and \texttt{push.\allowbreak{}py} refuses to start a second session while one is running or queued. A boot took about four minutes in the logged runs, with the model download taking 76 to 91 seconds (kaggle/\allowbreak{}logs/\allowbreak{}endpoint/\allowbreak{}). A pushed version that is stopped before it finishes does not become the notebook's default view on the Kaggle page; pinning it as the default version in the version history does (observed on the Kaggle page, 2026-11-15).

The n8n Cloud instance is a 14-day trial. After it ends the same workflow can be imported into a local n8n 1.x, which is what the launcher's local scripts do.

\subsection{Troubleshooting}

\begingroup\small
\begin{longtable}{>{\raggedright\arraybackslash}p{0.291\dimexpr\linewidth-6\tabcolsep\relax}>{\raggedright\arraybackslash}p{0.291\dimexpr\linewidth-6\tabcolsep\relax}>{\raggedright\arraybackslash}p{0.418\dimexpr\linewidth-6\tabcolsep\relax}}
\toprule
\textbf{symptom} & \textbf{cause found in practice} & \textbf{action} \\
\midrule
\endhead
\bottomrule
\endfoot
\endlastfoot
The workflow's model call fails right after a Kaggle restart & The tunnel address changed and the workflow still holds the old one & run \texttt{launcher\textbackslash{}start\textbackslash{}show\_\allowbreak{}model\_\allowbreak{}info.\allowbreak{}bat}; run \texttt{python tools\textbackslash{}golive\textbackslash{}set\_\allowbreak{}cloud\_\allowbreak{}qwen\_\allowbreak{}url.\allowbreak{}py} to post the new address \\
Kaggle status stays \texttt{QUEUED} for several minutes & the notebook is waiting for a GPU & wait; \texttt{push.\allowbreak{}py} will refuse a second push in the meantime \\
\texttt{watch\_\allowbreak{}deploy.\allowbreak{}py} prints \texttt{500 Server Error} from the Kaggle log stream & the Kaggle log API failed once & it reconnects by itself; check the status with \texttt{kaggle/\allowbreak{}tools/\allowbreak{}status.\allowbreak{}py} \\
Local n8n starts but the webhook does nothing & n8n 2.x does not register the webhook & install n8n 1.x (docs/\allowbreak{}guides/\allowbreak{}SETUP.md) \\
A fix pull request appears but CI on it is red & the model's fix parsed but was wrong & the retry loop tries again with thinking, up to the attempt cap; after that a comment says a person has to look \\
\bottomrule
\caption{Troubleshooting}\label{tab:10deploy-2}\\
\end{longtable}
\endgroup
